\documentclass[10pt,a4paper]{amsart}
\usepackage[margin=19mm]{geometry}
\usepackage{amsmath,amssymb,mathtools}
\usepackage[scr=rsfs,cal=euler]{mathalfa}
\usepackage{xfrac}
\usepackage[unicode,hidelinks,bookmarks=false]{hyperref}
\newcommand{\ie}{i.e.\ }
\newcommand{\cf}{cf.\ }
\newcommand{\resp}{respectively\ }
\newcommand{\Subsection}{\subsection}
\providecommand{\nobreakdash}{\nobreak\hbox{-}}
\setlength{\emergencystretch}{2em}
\multlinegap0pt
\usepackage{bm}
\usepackage{bbm}
\usepackage{dsfont}
\usepackage{xspace}
\usepackage{cancel}
\usepackage{mathtools}
\usepackage{amsthm}
\makeatletter
\@ifundefined{theorem}{\newtheorem{theorem}{Theorem}}{}
\@ifundefined{lemma}{\newtheorem{lemma}[theorem]{Lemma}}{}
\makeatother
\DeclareMathOperator{\Hom}{Hom}
\DeclareMathOperator{\Vol}{Vol}
\newcommand{\EE}{\mathbb{E}}
\newcommand{\PP}{\mathbb{P}}
\newcommand{\RR}{\mathbb{R}}
\newcommand{\ZZ}{\mathbb{Z}}
\newcommand{\calE}{\mathcal{E}}
\newcommand{\dd}{\,\mathrm{d}}
\title[Lipschitz Functions on Sparse Graphs II]{Lipschitz Functions on Sparse Graphs II}
\author{Samuel Korsky}
\date{Source: arXiv:2605.25515v1 (CC BY 4.0)}
\begin{document}
\maketitle

\noindent\emph{Source and licence.} Lipschitz Functions on Sparse Graphs II. Version arXiv:2605.25515v1 (CC BY 4.0), distributed under the Creative Commons Attribution 4.0 International licence, as recorded in the arXiv licence metadata for this exact version and shown on the abstract page https://arxiv.org/abs/2605.25515v1. Full rights notice: licenses/arxiv-2605.25515-CC-BY-4.0.txt.\par\smallskip

\noindent\textit{Editorial scope.} Counted unit: the Theorem whose source label is \texttt{thm:hypercube} in the article (source0.tex lines 894--900, source label \texttt{thm:hypercube}), delivered together with the article's own complete original proof of it (source0.tex lines 903--1054, one \texttt{proof} environment of 152 lines). No mathematics is written, repaired or re-derived: statement and proof are the article's own text line for line. Required context retained uncounted: lem:bounded-differences (lines 247--256); eq:rho-d (lines 274--277); lem:logistic (lines 295--305); lem:neighbor-extremes (lines 817--840). Every definition, display and label of those lines is retained unchanged. Omitted from this package and not redistributed: 1--246, 257--273, 278--294, 306--816, 841--893, 901--902, 1055--1106. Nothing inside a retained excerpt is omitted or altered: every retained body is delivered exactly as the article's lines are written, so no omission receipt of any kind accompanies this package. Citations: the retained excerpts carry no citation command at all, so no bibliography entry of the article is transcribed and no reference is redistributed. Reference adaptation: none was needed and none was made; every reference inside the delivered unit resolves inside it, so no unresolved reference or citation is produced. Printed slips kept literal: no printed word, symbol, hypothesis or number of the counted statement or of its proof was altered. Layout and class: the article's own class is \texttt{article} with packages including geometry, amsmath, amssymb, amsthm, mathtools, enumitem, hyperref, microtype; the wrapper is the lane's pinned \texttt{amsart} editorial wrapper at 10pt on A4, which restates the theorem-like environments the retained text needs and adds the article's own macro definitions that the retained text needs (\texttt{\textbackslash EE}, \texttt{\textbackslash Hom}, \texttt{\textbackslash PP}, \texttt{\textbackslash RR}, \texttt{\textbackslash Vol}, \texttt{\textbackslash ZZ}, \texttt{\textbackslash calE}, \texttt{\textbackslash dd}), with extra wrapper packages (bm, bbm, dsfont, xspace, cancel, mathtools). The delivered numbering is the unit's own. Assets: no retained span contains a figure environment, a graphics-inclusion command, a PDF-page inclusion, a table or a TikZ picture; no image file is copied into this package, the package declares no asset dependency and no image file is redistributed with it. Licence: version arXiv:2605.25515v1 of this article is distributed under the Creative Commons Attribution 4.0 International licence, as recorded in the arXiv licence metadata for this exact version and shown on the abstract page https://arxiv.org/abs/2605.25515v1. A full rights notice is delivered at licenses/arxiv-2605.25515-CC-BY-4.0.txt. Characters: every retained excerpt is pure ASCII, so the delivered engine, XeLaTeX with its default Latin Modern text font, reports no missing character.
\begin{lemma}[Bounded differences]\label{lem:bounded-differences}
Let $Y_1,\ldots,Y_m$ be independent random variables, and let $F=F(Y_1,\ldots,Y_m)$ be a real-valued function.  Suppose that changing only the $i$-th coordinate can change $F$ by at most $c_i$.  Then, for every $t>0$,
\[
  \PP(|F-\EE \left[F\right]|\ge t)\le 2\exp\left(-\frac{2t^2}{\sum_i c_i^2}\right).
\]
Moreover,
\[
  \log \EE \left[e^F\right]\le \EE \left[F\right]+\frac18\sum_i c_i^2.
\]
\end{lemma}

\begin{equation}\label{eq:rho-d}
  \rho_d(x):=\frac{1-e^{-d}}{(1+e^{-dx})(1+e^{d(x-1)})},
  \qquad x\in\RR.
\end{equation}

\begin{lemma}[Logistic boundary layers]\label{lem:logistic}
The density \eqref{eq:rho-d} satisfies
\begin{align}
  H(\rho_d)&=\frac{\pi^2}{3d}+O(d^{-2}),\label{eq:entropy-rho}\\
  Q(\rho_d)&=\frac{\pi^2}{3d^2}+O(d^{-3}).\label{eq:bad-rho}
\end{align}
Consequently,
\begin{equation}\label{eq:profile-gain}
  H(\rho_d)-\frac d2 \cdot Q(\rho_d)=\frac{\pi^2}{6d}+O(d^{-2}).
\end{equation}
\end{lemma}

\begin{lemma}[Neighbor extreme values]\label{lem:neighbor-extremes}
Let $T=T(d)\to\infty$ with $T^2=o(d)$, let
\[
  I=[-T/d,1+T/d],
\]
and let $\rho$ be the truncation and renormalization of $\rho_d$ to $I$.  Let $X_1,\ldots,X_d$ be independent samples from $\rho$, and set
\[
  M=\max_i X_i,
  \qquad
  m=\min_i X_i,
  \qquad
  R=m+1-M.
\]
Then
\[
  \EE \left[R\right]=o_d(d^{-1}),
  \qquad
  \EE \left[R^2\right]=O(d^{-2}),
\]
and consequently
\[
  \EE\left[\log(2-(M-m))\right]=\EE\left[\log(1+R)\right]\ge -o_d(d^{-1}).
\]
\end{lemma}

\begin{theorem}[Hypercube bounds]\label{thm:hypercube}
As $d\to\infty$,
\[
  \frac{\pi^2}{6d}+o(d^{-1})\le \log c(Q_d)
  \le \left(\frac34+o(1)\right)\frac{\log d}{d}.
\]
\end{theorem}

\begin{proof}
We first prove the lower bound.  Let $N=2^d$, and write the bipartition of $Q_d$ as $A\cup B$, with $|A|=|B|=N/2$.  By symmetry we may root later at a vertex in $A$.  Choose $T=T(d)\to\infty$ with $T^2=o(d)$, for instance $T=\log d$, and let $\rho$ be the truncation and renormalization of $\rho_d$ to
\[
  I=[-T/d,1+T/d].
\]
For $x\in I^A$ and $b\in B$, define
\[
  \ell_b(x):=\left|\bigcap_{a\sim b}[x_a-1,x_a+1]\right|.
\]
Since all values $x_a$ lie in an interval of length $1+2T/d<2$, the interval above is nonempty and
\[
  1-2T/d\le \ell_b(x)\le 2.
\]
Thus
\[
  Z_Q:=\int_{I^A}\prod_{b\in B}\ell_b(x)\dd x_A
\]
is an unrooted volume of valid Lipschitz assignments on $Q_d$: after choosing the values on $A$, each vertex $b\in B$ may be chosen independently in the interval of length $\ell_b(x)$.

\bigskip
\noindent
Let $X_a$, $a\in A$, be independent with density $\rho$, and put $W=-\log\rho$.  Then
\[
  Z_Q=\EE_\rho\left[\exp\left(\sum_{a\in A}W(X_a)\right)
  \prod_{b\in B}\ell_b(X)\right].
\]
Since $W$ is bounded by $O(T)$ under the truncated law, changing one variable $X_a$ changes $\sum_{a\in A}W(X_a)$ by at most $CT$.  Lemma~\ref{lem:bounded-differences} gives
\[
  \PP\left(\sum_{a\in A}W(X_a)-|A|H(\rho)\le -t\right)
  \le \exp\left(-\frac{2t^2}{(N/2)C^2T^2}\right).
\]
Taking $t=\gamma_d N/d$ with $\gamma_d\downarrow0$ and $\gamma_d^2N/(d^2T^2)\to\infty$ gives
\begin{equation}\label{eq:Q-W-conc}
  \sum_{a\in A}W(X_a)
  \ge |A|H(\rho)-o_d(N/d)
\end{equation}
with probability tending to one.

\bigskip
\noindent
We next control the product of interval lengths.  Fix $b\in B$, and let $M_b$ and $m_b$ be the maximum and minimum of the $d$ neighboring values $X_a$.  Then
\[
  \ell_b=2-(M_b-m_b).
\]
By Lemma~\ref{lem:neighbor-extremes},
\begin{equation}\label{eq:ell-expectation}
  \EE\left[\log\ell_b\right]\ge -o_d(d^{-1}).
\end{equation}
Now set
\[
  S(X):=\sum_{b\in B}\log\ell_b(X).
\]
Linearity of expectation and \eqref{eq:ell-expectation} give
\[
  \EE \left[S\right]\ge -o_d(N/d).
\]
Changing one coordinate $X_a$ affects only the $d$ terms with $b\sim a$.  Each term $\log\ell_b$ lies between $\log(1-2T/d)$ and $\log 2$, so changing one coordinate can change $S$ by at most $C d$.  Applying Lemma~\ref{lem:bounded-differences} with $|A|=N/2$ variables and $c_a=Cd$ gives, for every $t>0$,
\[
  \PP(S-\EE \left[S\right]\le -t)
  \le \exp\left(-\frac{2t^2}{(N/2)C^2d^2}\right).
\]
Choose any $\eta_d\downarrow0$ with $\eta_d^2N/d^4\to\infty$.  Taking $t=\eta_d N/d$ gives
\begin{equation}\label{eq:Q-length-conc}
  S(X)\ge -o_d(N/d)
\end{equation}
with probability tending to one.

\bigskip
\noindent
On the intersection $\calE$ of the high-probability events \eqref{eq:Q-W-conc} and \eqref{eq:Q-length-conc},
\[
  \exp\left(\sum_{a\in A}W(X_a)\right)
  \prod_{b\in B}\ell_b(X)
  \ge
  \exp\left(|A|H(\rho)-o_d(N/d)\right).
\]
Since $\PP_\rho(\calE)\to1$, its logarithm is $o_d(N/d)$.  Therefore
\[
  Z_Q\ge \PP_\rho(\calE)\exp\left(|A|H(\rho)-o_d(N/d)\right)
  =\exp\left(|A|H(\rho)-o_d(N/d)\right).
\]
Since $|A|=N/2$, $H(\rho)=H(\rho_d)+o_d(d^{-1})$, and Lemma~\ref{lem:logistic} gives $H(\rho_d)=\pi^2/(3d)+O(d^{-2})$, we obtain
\[
  Z_Q\ge \exp\left(N\left(\frac{\pi^2}{6d}+o(d^{-1})\right)\right).
\]
Slicing by the coordinate of a root in $A$ loses only a factor $|I|=1+o_d(1)$, so
\[
  \log c(Q_d)\ge\frac{\pi^2}{6d}+o(d^{-1}).
\]

\bigskip
\noindent
We prove the upper bound.  Let $N=2^d$, fix a root $o\in V(Q_d)$, and let $h$ be a positive integer.  Choose a fixed integer $L\ge5$ and set $M=Lh$.  Let $T_{M,h}$ be the graph with vertex set $\ZZ/M\ZZ$, with a loop at every vertex, and with an edge between two vertices when their cyclic distance is at most $h$.

\bigskip
\noindent
We first observe that rooted $h$-Lipschitz functions on $Q_d$ are in bijection with homomorphisms $\varphi:Q_d\to T_{M,h}$ satisfying $\varphi(o)=0$.  Reduction modulo $M$ gives one direction.  Conversely, given such a homomorphism, every oriented edge $uv$ has a unique lift increment $a_{uv}\in[-h,h]\cap\ZZ$ satisfying
\[
  \varphi(v)-\varphi(u)\equiv a_{uv}\pmod M,
\]
because $M>2h$.  Around every $4$-cycle, the sum of the four lifted increments is congruent to $0$ modulo $M$ and has absolute value at most $4h<M$, so it is actually $0$.  Since the cycle space of the hypercube is generated by its $4$-cycles, the lifted increments have zero sum around every cycle.  Integrating them from the root gives a unique rooted integer-valued $h$-Lipschitz function.

\bigskip
\noindent
The target $T_{M,h}$ is vertex-transitive.  Hence the number $N_{Q_d}(h)$ of rooted $h$-Lipschitz functions on $Q_d$ is
\begin{equation}\label{eq:Q-hom-rooted}
  N_{Q_d}(h)=\frac{\Hom(Q_d,T_{M,h})}{M}.
\end{equation}
By the Galvin--Tetali theorem, for every finite $d$-regular bipartite graph $G$ and every finite target graph $H$ with loops allowed,
\[
  \Hom(G,H)\le \Hom(K_{d,d},H)^{|V(G)|/(2d)}.
\]
Applying this with $G=Q_d$ and $H=T_{M,h}$ gives
\begin{equation}\label{eq:GT-Q}
  \Hom(Q_d,T_{M,h})
  \le \Hom(K_{d,d},T_{M,h})^{N/(2d)}.
\end{equation}
The same lifting argument applies to $K_{d,d}$, whose cycle space is also generated by $4$-cycles.  Therefore
\[
  \Hom(K_{d,d},T_{M,h})=M N_{K_{d,d}}(h)
  =L V_d h^{2d}+O_d(h^{2d-1}),
\]
where
\[
  V_d:=\Vol(P_{K_{d,d}})
\]
with one vertex of $K_{d,d}$ rooted.  Combining this with \eqref{eq:Q-hom-rooted} and \eqref{eq:GT-Q}, then taking the leading coefficient of $h^{N-1}$, gives
\begin{equation}\label{eq:Q-coeff-bound}
  c(Q_d)^{N-1}
  \le L^{-1}(L V_d)^{N/(2d)}.
\end{equation}
It remains to estimate $V_d$.

\bigskip
\noindent
Root one vertex in the left part of $K_{d,d}$.  If the values on the left part have range $r\le2$, then every vertex in the right part has an available interval of length $2-r$.  The volume of the set of left-part values with root fixed and range at most $r$ is $d r^{d-1}$; differentiating in $r$ gives
\begin{equation}\label{eq:Kdd-volume}
  V_d=d(d-1)\int_0^2 r^{d-2}(2-r)^d\,dr
  =2^{2d-1}d(d-1)B(d-1,d+1).
\end{equation}
By Stirling's formula,
\[
  V_d\sim \sqrt{\pi}\,d^{3/2}.
\]
Taking logarithms in \eqref{eq:Q-coeff-bound} and using $N=2^d$ yields
\[
  \log c(Q_d)
  \le \frac{\log V_d+O(1)}{2d}+o(1/d)
  =\left(\frac34+o(1)\right)\frac{\log d}{d}.
\]
This proves the theorem.
\end{proof}

\end{document}
